\section*{第 \ref{ch:g8:powers} 章——幂}

\begin{solution}{exo:g8:powers:1}
$2^6 = 64$; \quad $3^3 = 27$; \quad $10^5 = 100\,000$; \quad
$1^{100} = 1$; \quad $0.1^2 = 0.01$; \quad $6^0 = 1$.
\end{solution}

\begin{solution}{exo:g8:powers:2}
$(-3)^2 = 9$; \quad $-3^2 = -9$; \quad $(-1)^{15} = -1$; \quad
$(-5)^3 = -125$; \quad $\left(\frac23\right)^2 = \frac49$.
\end{solution}

\begin{solution}{exo:g8:powers:3}
$7^4 \times 7^5 = 7^9$; \quad
$\dfrac{2^9}{2^3} = 2^6$; \quad
$\left(10^3\right)^4 = 10^{12}$; \quad
$5^6 \times 5^{-2} = 5^4$; \quad
$3^4 \times 7^4 = 21^4$（\omterm{def:g8:powers:def}{指数}相同，底数相乘）。
\end{solution}

\begin{solution}{exo:g8:powers:4}
$10^{-2} = 0.01$; \quad
$4 \times 10^3 = 4000$; \quad
$2.5 \times 10^{-4} = 0.00025$; \quad
$10^0 = 1$.
\end{solution}

\begin{solution}{exo:g8:powers:5}
$10^5$; \quad $10^{-1}$; \quad $10^{10}$; \quad $10^{-6}$.
\end{solution}

\begin{solution}{exo:g8:powers:6}
$\dfrac{10^7 \times 10^{-3}}{10^2} = \dfrac{10^4}{10^2} = 10^2 =
100$.

$\dfrac{2^5 \times 2^4}{2^6} = \dfrac{2^9}{2^6} = 2^3 = 8$.

$\left(2^2\right)^3 \times 2^{-4} = 2^6 \times 2^{-4} = 2^2 = 4$.
\end{solution}

\begin{solution}{exo:g8:powers:7}
$2^3 \times 2^4 = 2^{12}$：\emph{错误}——\omterm{def:g8:powers:def}{指数}相加：
$2^7$.

$5^2 + 5^3 = 5^5$：\emph{错误}——和式没有对应的法则：
$25 + 125 = 150$，而 $5^5 = 3125$。

$(3^2)^4 = 3^8$：\emph{正确}。

$10^3 \times 10^3 = 100^3$：\emph{正确}——两者都等于 $10^6$
（左边：$10^{3+3}$；右边：$(10^2)^3$）。
\end{solution}

\begin{solution}{exo:g8:powers:8}
第 $4$ 天新增的人：第 3 天得知消息的 $3^3 = 27$ 人，
每人再告知三人，即新增 $3^4 = 81$ 人。截至第 4 天结束，知道消息的人共
$3 + 9 + 27 + 81 = 120$ 人。
\end{solution}

\begin{solution}{exo:g8:powers:9}
\emph{1.} 厚度：$2^{10} \times 10^{-4}$ m
$= 1024 \times 10^{-4}$ m $\approx 0.1$ m $= 10$ cm.

\emph{2.} $2^{42} \times 10^{-4} \approx 4.4 \times 10^{12} \times
10^{-4} = 4.4 \times 10^8$ m，大于 $3.8 \times 10^8$ m。经过
$42$ 次（纯属理论！）折叠，这叠纸就会超过月球。
\end{solution}

\begin{solution}{exo:g8:powers:10}
$2^{10} = 1024$；$10^3 = 1000$；$3^6 = 729$；$5^4 = 625$。从小到大排列为：
\[
5^4 < 3^6 < 10^3 < 2^{10} .
\]
\end{solution}

\begin{solution}{exo:g8:powers:11}
$2^{100} = \left(2^{10}\right)^{10} = 1024^{10}$，而
$10^{30} = \left(10^3\right)^{10} = 1000^{10}$。因为 $1024 > 1000$，
两边各取十次幂，不等号方向保持不变：
$2^{100} > 10^{30}$.
\end{solution}

\begin{solution}{pb:g8:powers:1}
\textbf{1.} 麦粒逐格翻倍，首格为
$1 = 2^0$ 粒：第 $k$ 格放 $2^{k-1}$ 粒
（见\cref{def:g8:powers:def}）。第 $64$ 格放 $2^{63}$ 粒。

\textbf{2.} $S_1 = 1$，$S_2 = 1 + 2 = 3$，$S_3 = 3 + 4 = 7$，
$S_4 = 7 + 8 = 15$，$S_5 = 15 + 16 = 31$：总比
下一个 $2$ 的幂小 1（$2$、$4$、$8$、$16$、$32$）。猜想：
$S_n = 2^n - 1$.

\textbf{3.} 将 $S_n$ 的每一项都乘以 2，每个幂的\omterm{def:g8:powers:def}{指数}都增加
1（$2 \times 2^{k} = 2^{k+1}$）：
\begin{align*}
S_n &= 1 + 2 + 4 + \dots + 2^{n-1}, \\
2 \times S_n &= \phantom{1 + {}} 2 + 4 + \dots + 2^{n-1} + 2^n .
\end{align*}
将第二式减去第一式，从 $2$
到 $2^{n-1}$ 的每一项在两式中都出现，全部相消：
\[
2 \times S_n - S_n = 2^n - 1,
\qquad\text{即}\qquad
S_n = 2^n - 1 .
\]

\textbf{4.} 棋盘上麦粒的总数为 $S_{64} = 2^{64} - 1$ 粒。
设想第 65 格放 $2^{64}$ 粒，则整个棋盘
的粒数比那一格所应放的恰好少 1 粒。

\textbf{5.} 前\omterm{def:g3:division:half}{一半}棋盘共有 $S_{32} = 2^{32} - 1$ 粒。
整个棋盘共有 $2^{64} - 1$ 粒，因此后\omterm{def:g3:division:half}{一半}共有
\[
\left(2^{64} - 1\right) - \left(2^{32} - 1\right)
= 2^{64} - 2^{32}
= 2^{32} \times \left(2^{32} - 1\right)
\]
（提取公因式 $2^{32}$，利用 $2^{32} \times 2^{32} = 2^{64}$，
\cref{thm:g8:powers:rules}）：恰好是前\omterm{def:g3:division:half}{一半}的 $2^{32}$
倍。

\textbf{6.} 由\omterm{def:g8:powers:def}{指数}法则，$2^{64} = 2^{4 + 60} =
2^4 \times \left(2^{10}\right)^6$。因为 $2^{10} = 1024 > 10^3$，于是有
\[
2^{64} > 16 \times \left(10^3\right)^6 = 16 \times 10^{18}
= 1.6 \times 10^{19} .
\]

\textbf{7.} 超过 $1.6 \times 10^{19}$ 粒，每粒
重 $5 \times 10^{-2}$ g：
\[
1.6 \times 10^{19} \times 5 \times 10^{-2}
= 8 \times 10^{17} \text{ g} .
\]
除以每吨 $10^6$ g，得超过 $8 \times 10^{11}$
吨——即八千亿吨。

\textbf{8.} $\dfrac{8 \times 10^{11}}{8 \times 10^{8}} = 10^3$:
国王许诺至少给出\emph{千年}的当今
全世界的全部收成。（传说中，他的谋士们
也曾这样告诉过他。）

\textbf{9.} $2^{19} = 2^9 \times 2^{10} = 512 \times 1024 =
524\,288 < 10^6$，而 $2^{20} = \left(2^{10}\right)^2 =
1024^2 = 1\,048\,576 > 10^6$。于是最小的\omterm{def:g8:powers:def}{指数}为
$n = 20$：二十次翻倍便超过了一百万。

\textbf{10.} 第 $n$ 天的报酬为 $2^{n-1}$ 分（第 1 天：
$2^0 = 1$）。现在
\[
2^{26} = 2^6 \times \left(2^{10}\right)^2
= 64 \times 1\,048\,576 = 67\,108\,864 < 10^8,
\]
\[
2^{27} = 2 \times 2^{26} = 134\,217\,728 > 10^8 .
\]
因此，每日报酬首次超过 $10^8$ 分，是在 $n - 1 = 27$ 时：
即第 $28$ 天。

\textbf{11.} $21 = 16 + 4 + 1$：砝码分别为 $16$、$4$ 和 $1$ g。
$27 = 16 + 8 + 2 + 1$：砝码分别为 $16$、$8$、$2$ 和 $1$ g。

\textbf{12.} 由第 3 题，砝码 $1, 2, 4, 8$ 合计
重 $S_4 = 2^4 - 1 = 15$ g。因此若没有 $16$ g 的砝码，
商人无法称出超过 $15$ g 的重量：任何 $16$ g 或更大的目标
都必须用到它。而不超过 $15$ g 的目标则不能用它，因为
$16$ g 砝码本身就已超过目标。因此最大砝码的选取
是\emph{必然的}。剩下的问题是
不超过 $15$ g 的目标，须由 $1, 2, 4, 8$ 凑成——而
同一论证依次重复：$8$ 的选取同样必然（若剩余
目标 $\geq 8$ 则用，否则不用，因为 $1 + 2 + 4 = 7$），
然后是 $4$（因为 $1 + 2 = 3$），再是 $2$，最后是 $1$。

\textbf{13.} 按照这些必然选择，
每一步之后剩下的目标至多等于剩下砝码的总和，
因此过程以\omterm{def:g3:division:remainder}{余数} $0$ 结束：从 $1$ 到
$31$ 的每个目标都能达到。又因为沿途每一步都是必然选择，
其他任何砝码组合都不能达到同一目标：
每个从 $1$ 到 $31$ 的数，都可写成互不相同的
$2$ 的幂之和，且这种写法存在并\emph{唯一}。

\textbf{14.} 六个砝码共重 $S_6 = 2^6 - 1 = 63$ g，
同样的必然选择论证适用于从 $1$ 到
$63$ g 的每一个目标。对于 $45$：目标 $\geq 32$，故用 $32$；剩余
$13 < 16$，故不用 $16$；用 $8$（剩余 $5$），用 $4$（剩余
$1$），不用 $2$，用 $1$：
\[
45 = 32 + 8 + 4 + 1 = 2^5 + 2^3 + 2^2 + 2^0 .
\]

\textbf{15.} 每个方格选取 $0$ 或 $1$，共 $64$ 个方格，
相当于选取哪些幂 $2^0, 2^1, \dots, 2^{63}$ 参与求和
——这恰是商人用 $64$
个砝码称重的过程。可存储的最\omterm{def:g6:decimals:places}{小数}是 $0$（所有方格均为
$0$）；最大的数是\emph{所有}幂之和，即
国王的总数 $S_{64} = 2^{64} - 1$（第一部分）。由
必然选择的论证，其间每个整数都被取到
恰好一次：一台 $64$ 位机器正好能存储
从 $0$ 到 $2^{64} - 1$ 的全部整数。
\end{solution}
